\subsection{Negative Integer Exponents} 
Next consider power functions of the form
\[
f(x)=a(x-h)^{-n}+k=\makeblank[2.5in]
\]
where $h,k$ are \makeblank[1in] numbers and $n$ is a \makeblank[1in] number. (This means that $-n$ is a \makeblank.)
\subsubsection*{Domain}
This is a \makeblank[1in] function, and its \makeblank is \makeblanksmall only when $x=\makeblanksmall$. So its domain is $\left\{x\suchthat\makeblank[1in]\right\}$ (or \makeblank[1.5in] in interval notation).
\subsubsection*{Range}
\begin{itemize}
\item \makeblank[1in] is never \makeblanksmall, so \makeblank is not in the range of $f$.
\item If $n$ is \makeblank[1in], $(x-h)^n$ can be \makeblank[1.5in], so the range of $f$ is $\left\{x\suchthat\makeblank[1in]\right\}$ (or \makeblank[1.5in] in interval notation).
\item If $n$ is \makeblank[1in], $(x-h)^n$ must be \makeblank[1.5in], so the range of $f$ depends on $a$ and $k$.
\begin{itemize}
\item If $a$ is \makeblank[1in], the range of $f$ is \makeblank.
\item If $a$ is \makeblank[1in], the range of $f$ is \makeblank.
\end{itemize}
\end{itemize}
\subsubsection*{Inverse}
If $n$ is \makeblank[1in], $f(x)=x^{-n}$ has an inverse.
\begin{lpic}{}{3}
\end{lpic}
If $n$ is \makeblank[1in], we have to split up the domain of $f(x)=x^{-n}$.
\begin{lpic}{}{4}
\twocol{-1}{4}{$x>0$}{$x<0$}
\end{lpic}
